\documentclass[11pt]{article}

\usepackage[margin=1in]{geometry}
\usepackage{amsmath}
\usepackage{amssymb}
\usepackage{booktabs}
\usepackage[hidelinks]{hyperref}

\title{Bioinformatics of Exoprotocells}
\author{LIFEORIG working note}
\date{\today}

\begin{document}

\maketitle

\section{Purpose}

This note connects the Titan protocell framework to a bioinformatics program. The existing Titan notes define a methane-pool environment, a reaction network, phase allocation, dense organic phases, and interface-stabilized protocell candidates. The next computational question is how to analyze the resulting population once many candidate compartments exist.

The central proposal is that an exoprotocell does not need a DNA or RNA genome in order to generate an informatic state at the population level. The protocell itself is first a dynamical physical-chemical object. Its state evolves by internal reaction kinetics, transport across the boundary, structural change of the compartment, and coupling to the local environment. A useful abstract state is:
\begin{equation}
\mathcal{P}_k =
\left(
\mathbf{C}^{P_k},
\mathbf{u}^{P_k},
\mathbf{z}^{P_k},
\mathbf{E}(z_k,t),
\mathcal{H}_k
\right),
\end{equation}
where $\mathbf{C}^{P_k}$ is the internal chemical composition, $\mathbf{u}^{P_k}$ is the boundary/interface structure and permeability state, $\mathbf{z}^{P_k}$ contains physical object variables such as radius, volume, and surface area, $\mathbf{E}(z_k,t)$ is the local external environment at the protocell position, and $\mathcal{H}_k$ is the object history.

The active part of the chemical network is not assumed to be fixed inside the protocell. It is an emergent subnetwork selected by the current state:
\begin{equation}
\mathcal{R}_{\rm active}^{P_k}(t)
=
\mathcal{A}
\left(
\mathbf{C}^{P_k}(t),
\mathbf{u}^{P_k}(t),
\mathbf{z}^{P_k}(t),
\mathbf{E}(z_k,t)
\right).
\end{equation}
Bioinformatics enters after this dynamical step: it analyzes the distribution of protocell states, structures, active subnetworks, and lineages across a population.

\section{Relation to the Titan Theory}

The Titan design note separates chemical identity from physical phase allocation:
\begin{equation}
\text{chemical species classification}
\neq
\text{physical phase allocation}.
\end{equation}
This distinction is also the starting point for exoprotocell bioinformatics. The relevant unit is not only a molecule sequence. It is a coupled object:
\begin{equation}
\text{species}
+
\text{phase}
+
\text{reaction role}
+
\text{retention}
+
\text{history}.
\end{equation}

The Titan model promotes compartments by local threshold rules:
\begin{align}
C_{\rm coa}(z_j,t) &\ge C_{\rm coa}^{\rm crit}
\quad \Rightarrow \quad
\text{dense organic object},\\
C_I(z_j,t) &\ge C_I^{\rm crit}
\quad \Rightarrow \quad
\text{interface-stabilized protocell candidate}.
\end{align}
Bioinformatics begins after this promotion step. Each promoted object becomes a sample in a population dataset.

\section{Dynamical State and Informatic State}

For an ordinary cell, bioinformatics often starts from sequence data. For a Titan exoprotocell candidate, the first object is instead the dynamical state:
\begin{equation}
\frac{d\mathbf{C}^{P_k}}{dt}
=
\mathbf{S}\mathbf{r}^{P_k}
+
\mathbf{J}_{\rm in}^{P_k}
-
\mathbf{J}_{\rm out}^{P_k}
-
\mathbf{D}^{P_k},
\end{equation}
where $\mathbf{S}$ is the stoichiometric matrix, $\mathbf{r}^{P_k}$ is the vector of internal reaction rates, $\mathbf{J}_{\rm in}^{P_k}$ and $\mathbf{J}_{\rm out}^{P_k}$ are exchange fluxes with the local environment, and $\mathbf{D}^{P_k}$ represents degradation, leakage, burial, or photochemical damage.

The exchange terms depend on both the protocell structure and the local environment:
\begin{align}
J_{{\rm in},i}^{P_k}
&=
\frac{A_k}{V_k}
P_i
\left(\mathbf{u}^{P_k},\mathbf{E}(z_k,t)\right)
C_i^{\rm ext}(z_k,t),\\
J_{{\rm out},i}^{P_k}
&=
\frac{A_k}{V_k}
P_i
\left(\mathbf{u}^{P_k},\mathbf{E}(z_k,t)\right)
C_i^{P_k}(t).
\end{align}
Thus two protocells with the same internal composition can diverge if they have different shell structures, surface coverage, permeability profiles, birth positions, or local chemical environments.

The informatic state is not an additional physical state inside the protocell. It is a feature representation extracted from the dynamical state and its history:
\begin{equation}
\mathbf{x}_k =
\left[
\mathbf{c}_k,
\mathbf{p}_k,
\mathbf{a}_k,
\mathbf{g}_k,
\mathbf{s}_k,
\mathbf{h}_k
\right].
\end{equation}

\begin{center}
\begin{tabular}{p{0.16\linewidth}p{0.70\linewidth}}
\toprule
Block & Meaning \\
\midrule
$\mathbf{c}_k$ & internal composition vector for protocell candidate $k$ \\
$\mathbf{p}_k$ & phase and localization vector: internal, interface, adsorbed, dense, external-nearby \\
$\mathbf{a}_k$ & activity vector: active reactions, fluxes, catalytic factors, exchange rates \\
$\mathbf{g}_k$ & graph features of the currently active subnetwork \\
$\mathbf{s}_k$ & structural vector: shell composition, permeability, radius, surface area, retention times \\
$\mathbf{h}_k$ & history features: birth position, parent object, exposure, division, fusion, loss \\
\bottomrule
\end{tabular}
\end{center}

The object to track is therefore not a private genome for each protocell. It is an effective organization that constrains which reactions from the shared network are expressed by the protocell state:
\begin{equation}
O_k^{\rm eff}(t) =
\left(
\mathcal{R}_{\rm active}^{P_k},
\mathcal{C}_{\rm catalysts}^{P_k},
\mathbf{P}^{P_k},
\mathbf{u}^{P_k},
\mathbf{E}(z_k,t)
\right).
\end{equation}
Here $\mathbf{P}^{P_k}$ is the permeability profile. The boundary is included because, for non-aqueous protocells, chemical access and retention are part of the effective inherited organization.

\section{Polymer-Generated Catalytic Diversity}

The simulation has one global reaction network. Let
\begin{equation}
\mathcal{R}=\{r_1,\ldots,r_M\},
\qquad
\mathcal{S}=\{s_1,\ldots,s_N\}.
\end{equation}
Every protocell uses the same $\mathcal{R}$ and $\mathcal{S}$. Diversity is not introduced by giving each protocell a different network. Diversity is introduced by allowing different protocells to develop different internal polymer populations.

Let the monomer alphabet be
\begin{equation}
\mathcal{M}=\{m_1,\ldots,m_A\}.
\end{equation}
A polymer sequence is
\begin{equation}
\pi_a=(\sigma_1,\ldots,\sigma_{L_a}),
\qquad
\sigma_j\in\mathcal{M},
\end{equation}
where $L_a$ is the chain length. The number of possible sequences of length $L$ is
\begin{equation}
|\Pi_L|=A^L.
\end{equation}
Thus longer polymers create a combinatorial space of possible structures. Most sequences may be inactive, but some sequences can fold, aggregate, adsorb, or expose chemical groups in ways that alter reaction rates.

For protocell $k$, define the internal polymer copy number:
\begin{equation}
N^{\rm poly}_{a,k}(t)
=
\text{number of copies of polymer sequence }\pi_a
\text{ inside protocell }k.
\end{equation}
The polymer population vector is
\begin{equation}
\mathbf{N}^{\rm poly}_k(t)
=
\left(
N^{\rm poly}_{1,k},
\ldots,
N^{\rm poly}_{N_{\rm poly},k}
\right).
\end{equation}

The key biochemical object is the polymer-to-reaction catalytic coefficient:
\begin{equation}
\Gamma_{ar}
=
\text{catalytic strength of polymer }\pi_a
\text{ on reaction }r.
\end{equation}
This coefficient can be zero for most polymer-reaction pairs. It can depend on sequence length, composition, conformation, phase localization, and solvent compatibility. In a first model it can be assigned statistically:
\begin{equation}
\Gamma_{ar}\sim P_\Gamma(L_a,\chi_a,r),
\end{equation}
where $\chi_a$ is a sequence descriptor such as nitrile fraction, hydrophobic fraction, charge-like polarity class, or HCN-polymer motif count.

The effective catalytic activity of reaction $r$ inside protocell $k$ is then
\begin{equation}
\Gamma^{\rm eff}_{r,k}(t)
=
\Gamma^0_r
+
\sum_a
\Gamma_{ar}
\frac{N^{\rm poly}_{a,k}(t)}{V_k(t)}
\,\theta_{a,k}(t),
\end{equation}
where $\Gamma^0_r$ is the background uncatalyzed or externally catalyzed contribution, $V_k$ is protocell volume, and $\theta_{a,k}$ is the fraction of polymer $\pi_a$ in an active location or conformation.

The reaction rate inside protocell $k$ can then be written as
\begin{equation}
R_{r,k}(t)
=
\Gamma^{\rm eff}_{r,k}(t)
F_r\!\left(\mathbf{E}(z_k,t)\right)
\prod_i
\left[
C^{P_k}_i(t)
\right]^{\nu_{ir}},
\end{equation}
where $F_r(\mathbf{E})$ is the extrinsic environmental factor and $\Gamma^{\rm eff}_{r,k}$ is the intrinsic protocell catalytic factor. This split is essential:
\begin{equation}
\text{environmental forcing}
\neq
\text{intrinsic protocell diversity}.
\end{equation}

A protocell-specific active reaction set can now be defined without changing the global network:
\begin{equation}
\mathcal{R}^{P_k}_{\rm active}(t)
=
\{r\in\mathcal{R}: R_{r,k}(t)>R_{\rm min}\}.
\end{equation}
The active set differs because $\mathbf{N}^{\rm poly}_k$, $\theta_{a,k}$, and $\Gamma^{\rm eff}_{r,k}$ differ, not because $\mathcal{R}$ differs.

\section{Measuring Protocell Diversity}

The intrinsic diversity variable should be
\begin{equation}
X_k(t)
=
\left(
\mathbf{N}^{\rm poly}_k(t),
\boldsymbol{\Gamma}^{\rm eff}_k(t),
\mathbf{C}^{P_k}(t),
\mathbf{P}_k(t),
\boldsymbol{\theta}_k(t)
\right),
\end{equation}
where
\begin{align}
\boldsymbol{\Gamma}^{\rm eff}_k(t)
&=
\left(\Gamma^{\rm eff}_{1,k},\ldots,\Gamma^{\rm eff}_{M,k}\right),\\
\mathbf{P}_k(t)
&=
\left(P_{1,k},\ldots,P_{N,k}\right),\\
\boldsymbol{\theta}_k(t)
&=
\left(\theta_{1,k},\ldots,\theta_{N_{\rm poly},k}\right).
\end{align}
Here $\mathbf{P}_k$ is the permeability vector for species transport across the protocell boundary. The vector $\boldsymbol{\theta}_k$ records which polymers are actually folded, surface-bound, dense-phase localized, or otherwise catalytically available.

A concrete distance between protocells is:
\begin{align}
d_{\rm proto}(k,l)
=&
w_{\rm poly}
\frac{
\|\mathbf{N}^{\rm poly}_k-\mathbf{N}^{\rm poly}_l\|_1
}{
N^{\rm poly}_k+N^{\rm poly}_l+\epsilon
}
\nonumber\\
&+
w_\Gamma
\|\boldsymbol{\Gamma}^{\rm eff}_k-\boldsymbol{\Gamma}^{\rm eff}_l\|_2
+
w_C
\|\mathbf{C}^{P_k}-\mathbf{C}^{P_l}\|_1
\nonumber\\
&+
w_P
\|\mathbf{P}_k-\mathbf{P}_l\|_2
+
w_\theta
\|\boldsymbol{\theta}_k-\boldsymbol{\theta}_l\|_1.
\end{align}
This is the object that should be clustered in a bioinformatics analysis. It measures intrinsic protocell differences. Local environment can be stored separately as metadata:
\begin{equation}
\mathbf{E}_k(t)=\mathbf{E}(z_k,t).
\end{equation}
The environment may change measured activity, but it should not be confused with protocell identity.

\section{Exoprotocell Quasispecies}

The population-level object is the quasispecies. In this framework, a quasispecies is not a cloud of nucleic-acid sequences. It is a cloud of protocell states whose main intrinsic diversity comes from retained polymer populations and the catalytic matrices generated by those polymers:
\begin{equation}
\mathcal{Q}(t)
=
\left\{
\left(
\mathbf{N}^{\rm poly}_k,
\boldsymbol{\Gamma}^{\rm eff}_k,
\mathbf{C}^{P_k},
\mathbf{P}_k,
\boldsymbol{\theta}_k
\right)
\right\}_{k=1}^{N_P(t)}.
\end{equation}

The local environment remains important, but it is not the definition of protocell type. A protocell at the UV-exposed surface and a protocell attached to a mineral surface may have different observed rates because $\mathbf{E}(z_k,t)$ is different. That is environmental heterogeneity. Protocell diversity requires differences in $X_k(t)$ itself: different retained polymers, different polymer conformational availability, different boundary permeability, and therefore different $\boldsymbol{\Gamma}^{\rm eff}_k$.

Let $\phi_\alpha(t)$ be a coarse-grained protocell type, obtained by clustering records $\mathbf{x}_k$ into an equivalence class:
\begin{align}
\phi_\alpha
=
\big[
&\text{polymer-population fingerprint},
\text{effective-catalysis fingerprint},
\nonumber\\
&\text{internal-composition fingerprint},
\text{permeability fingerprint},
\text{active-conformation fingerprint}
\big]_\alpha.
\end{align}
The quasispecies distribution is then
\begin{equation}
q_\alpha(t)
=
\frac{N_\alpha(t)}{\sum_\beta N_\beta(t)}.
\end{equation}

A useful first dynamical form is:
\begin{equation}
\frac{dq_\alpha}{dt}
=
\sum_\beta W_{\alpha\beta}q_\beta
-
\bar{W}q_\alpha,
\end{equation}
where $W_{\alpha\beta}$ is an effective transition-selection operator. It includes polymer growth, polymer cleavage, polymer loss, sequence-dependent catalysis, protocell growth, fragmentation, fusion, and boundary modification. Environmental relocation can change the observed rates, but a transition between intrinsic types requires a change in $X_k$. The mean population persistence is:
\begin{equation}
\bar{W}
=
\sum_{\alpha,\beta} W_{\alpha\beta}q_\beta.
\end{equation}

In this sense, mutation is generalized. A protocell type can change because its polymer population changes:
\begin{equation}
\mathbf{N}^{\rm poly}_k
\rightarrow
\mathbf{N}^{\rm poly}_k+\Delta\mathbf{N}^{\rm poly}_k,
\end{equation}
which induces
\begin{equation}
\boldsymbol{\Gamma}^{\rm eff}_k
\rightarrow
\boldsymbol{\Gamma}^{\rm eff}_k+\Delta\boldsymbol{\Gamma}^{\rm eff}_k.
\end{equation}
The formal reaction network remains fixed. The quasispecies is the population distribution over polymer-generated catalytic states.

\section{Exoprotocell Omics}

The following analogs are useful for implementation and analysis.

\begin{center}
\small
\begin{tabular}{p{0.24\linewidth}p{0.55\linewidth}p{0.13\linewidth}}
\toprule
Earth bioinformatics term & Exoprotocell analog & Data object \\
\midrule
Genome-like state & retained polymer population inside the protocell & vector \\
Transcriptome-like state & reaction-rate vector induced by the polymer population & vector \\
Proteome-like state & catalytically active polymers, conformations, and surface sites & set/vector \\
Metabolome & internal chemical composition & vector \\
Lipidome & interface or shell composition & vector \\
Microbiome & local population of interacting protocells & graph \\
Phylogeny & lineage of birth, fusion, fragmentation, inheritance & tree/DAG \\
\bottomrule
\end{tabular}
\end{center}

These labels should be used carefully. They are computational analogies, not claims that Titan protocells possess terrestrial biology.

\section{Population Bioinformatics}

A simulation output should be convertible into a population table:
\begin{equation}
\mathcal{D}(t)
=
\{ \mathbf{x}_1(t),\ldots,\mathbf{x}_{N_P(t)}(t) \}.
\end{equation}
The minimal columns should be:
\begin{equation}
\begin{aligned}
\texttt{id},\quad
\texttt{parent\_id},\quad
\texttt{birth\_time},\quad
\texttt{birth\_z},\quad
\texttt{radius},\quad
\texttt{interface\_type},\\
\texttt{polymer\_fingerprint},\quad
\texttt{gamma\_eff\_fingerprint},\\
\texttt{permeability\_fingerprint},\quad
\texttt{retention\_score}.
\end{aligned}
\end{equation}

This enables standard population analyses:
\begin{itemize}
\item clustering of protocell states,
\item trajectory analysis through chemical state space,
\item detection of recurrent polymer-generated catalytic motifs,
\item lineage reconstruction after fragmentation or division-like events,
\item survival analysis of compartment lifetimes,
\item association between interface composition and persistence.
\end{itemize}

\section{Lineages Without Genes}

In early exoprotocell simulations, inheritance may occur through partial material continuity rather than sequence copying. If protocell $k$ fragments into daughters $a$ and $b$, inheritance can be measured by composition overlap:
\begin{equation}
H_{a|k}
=
\frac{\sum_i \min(C_i^{P_a},C_i^{P_k})}
{\sum_i C_i^{P_a}}.
\end{equation}
Polymer inheritance can be measured by polymer population overlap:
\begin{equation}
J_{\rm poly}(a,k)
=
\frac{
\sum_p \min(N^{\rm poly}_{p,a},N^{\rm poly}_{p,k})
}{
\sum_p \max(N^{\rm poly}_{p,a},N^{\rm poly}_{p,k})
}.
\end{equation}
Effective catalytic inheritance can be measured by the similarity of the induced catalytic matrices:
\begin{equation}
J_{\Gamma}(a,k)
=
1-
\frac{
\|\boldsymbol{\Gamma}^{\rm eff}_a-\boldsymbol{\Gamma}^{\rm eff}_k\|_2
}{
\|\boldsymbol{\Gamma}^{\rm eff}_a\|_2+
\|\boldsymbol{\Gamma}^{\rm eff}_k\|_2+\epsilon
}.
\end{equation}
The lineage graph is then a directed acyclic graph:
\begin{equation}
\mathcal{L}=(V,E),
\end{equation}
where vertices are protocell objects and edges record division, fragmentation, fusion, or persistent material transfer.

\section{Feature Engineering for Titan}

For Titan specifically, the most important features are not fast growth features. They are persistence features:
\begin{align}
S_{{\rm persist},k}
&=
w_1\tau_{{\rm survival},k}
+w_2\tau_{{\rm retention},k}
+w_3\theta_{I,k}
+w_4M_{{\rm polymer},k}
\nonumber\\
&\quad
+w_5R_{{\rm catalytic},k}
-w_6D_{{\rm photochemical},k}
-w_7L_{{\rm sediment},k}.
\end{align}

The corresponding bioinformatics features should include:
\begin{itemize}
\item nitrile-rich composition fraction,
\item HCN-polymer or tholin-like polymer fraction,
\item interface-forming species fraction,
\item polymer sequence-length distribution,
\item polymer sequence diversity or Shannon entropy,
\item effective catalytic matrix $\boldsymbol{\Gamma}^{\rm eff}_k$,
\item catalytic closure score induced by retained polymers,
\item RAF or autocatalytic subset size,
\item retention time distribution over key species,
\item phase entropy, measuring whether material is concentrated or dispersed,
\item environmental birth context: surface, sediment, adsorbed layer, or dense organic phase.
\end{itemize}

One useful phase entropy is
\begin{equation}
H_{{\rm phase},k}
=
-\sum_m q_{m,k}\log q_{m,k},
\end{equation}
where $q_{m,k}$ is the fraction of protocell-associated material in phase $m$. Low entropy may indicate strong compartmentalization; high entropy may indicate diffuse material with weak protocell identity.

\section{Implementation Proposal}

The implementation should not replace the chemistry or dynamics layer. It should be an analysis layer that consumes simulation state snapshots.

\subsection{New data objects}

Recommended first objects:
\begin{align}
&\texttt{ExoprotocellRecord},\quad
\texttt{ExoprotocellPopulation},\quad
\texttt{PolymerPopulation},
\nonumber\\
&\texttt{EffectiveCatalyticMatrix},\quad
\texttt{LineageGraph}.
\end{align}

\begin{center}
\small
\begin{tabular}{p{0.30\linewidth}p{0.58\linewidth}}
\toprule
Object & Responsibility \\
\midrule
\texttt{ExoprotocellRecord} & store one protocell candidate snapshot \\
\texttt{ExoprotocellPopulation} & store all protocell snapshots at time $t$ \\
\texttt{PolymerPopulation} & store polymer sequence counts $N^{\rm poly}_{a,k}$ for one protocell \\
\texttt{EffectiveCatalyticMatrix} & compute $\boldsymbol{\Gamma}^{\rm eff}_k$ from retained polymers \\
\texttt{LineageGraph} & record parent, daughter, fusion, and transfer relations \\
\bottomrule
\end{tabular}
\end{center}

\subsection{Analysis functions}

Recommended first functions:
\begin{equation}
\begin{gathered}
\texttt{fingerprint\_composition(record)},\\
\texttt{fingerprint\_polymer\_population(record)},\\
\texttt{compute\_effective\_catalytic\_matrix(record)},\\
\texttt{compute\_closure\_score(record)},\\
\texttt{compute\_retention\_score(record)},\\
\texttt{compare\_records(record\_a, record\_b)},\\
\texttt{build\_lineage(records)}.
\end{gathered}
\end{equation}

The output should be plain tables and sparse matrices first. More elaborate visualization can come later.

\section{Connection to Current LIFEORIG Code}

The existing code already has several useful anchors:
\begin{itemize}
\item \texttt{MoleculeModel}, \texttt{BinaryPolymer}, \texttt{MultiPolymer}, and \texttt{ReferenceMolecule} can define chemical identity records.
\item \texttt{ChemicalNetworkParser} can provide parsed species and reaction records.
\item ACF/RAF-style network routines can provide closure and autocatalytic subset features.
\item Titan environment parsers provide the environmental context needed for birth and persistence features.
\item The future \texttt{TitanPhaseAllocator} can provide phase labels and phase inventories.
\end{itemize}

The key design rule is:
\begin{equation}
\text{bioinformatics layer}
=
\text{read-only analysis of simulated protocell states}.
\end{equation}
It should not own the chemistry, the environmental dynamics, or the phase transition rules.

\section{Minimal First Milestone}

The first useful milestone is a post-processing pipeline:
\begin{enumerate}
\item Run a Titan methane-pool simulation.
\item Save protocell candidate snapshots at selected times.
\item Convert each candidate into an \texttt{ExoprotocellRecord}.
\item Build polymer-population fingerprints and effective-catalytic-matrix fingerprints.
\item Cluster candidates by $d_{\rm proto}(k,l)$.
\item Build a lineage graph from parent-child or fragmentation metadata.
\item Report persistence-associated features.
\end{enumerate}

The scientific output is not a claim of life. It is a map of which chemical organizations persist, recur, and transmit structure under Titan-like constraints.

\section{References and Local Inputs}

\subsection*{Local LIFEORIG notes and code}
\begin{itemize}
\item \texttt{doc/TITAN/titan\_methane\_pool\_protocell\_design.tex}
\item \texttt{doc/TITAN/titan\_catalysis\_framework\_latex.pdf}
\item \texttt{doc/TITAN/titan\_protocell\_framework.pdf}
\item \texttt{src/chemical\_types/molecule\_model.py}
\item \texttt{src/input\_data/chemical\_network\_parser.py}
\item \texttt{src/RAF/RAF\_engine.py}
\end{itemize}

\subsection*{Literature for polymer-generated protocell diversity}

\begin{itemize}
\item Walker, S. I., Grover, M. A., and Hud, N. V. Universal sequence replication, reversible polymerization and early functional biopolymers: A model for the initiation of prebiotic sequence evolution. 2012. \url{https://arxiv.org/abs/1203.4625}
\item Guseva, E. A., Zuckermann, R. N., and Dill, K. A. How did prebiotic polymers become informational foldamers? 2016. \url{https://arxiv.org/abs/1604.08890}
\item Tkachenko, A. V. and Maslov, S. Onset of natural selection in auto-catalytic heteropolymers. 2017. \url{https://arxiv.org/abs/1710.06385}
\item Hordijk, W. and Steel, M. Autocatalytic sets in polymer networks with variable catalysis distributions. 2016. \url{https://arxiv.org/abs/1605.03919}
\item Kocher, C. D. and Dill, K. A. The prebiotic emergence of biological evolution. 2023. \url{https://arxiv.org/abs/2311.13650}
\item Segr\'e, D., Ben-Eli, D., and Lancet, D. Compositional genomes: Prebiotic information transfer in mutually catalytic noncovalent assemblies. \textit{Proceedings of the National Academy of Sciences}, 2000.
\item Eigen, M. and Schuster, P. The hypercycle: A principle of natural self-organization. 1979.
\item Eigen, M., McCaskill, J., and Schuster, P. Molecular quasi-species. \textit{Journal of Physical Chemistry}, 1988.
\end{itemize}

\end{document}
